% Source: physical PDF pp. 276--278 (printed pp. 253--255). Coverage: complete Section 20.1, from its heading through immediately before Section 20.2; Eqs. (20.1.1)--(20.1.9), citation [1], one footnote, no figures. Uncertainty: none.
\subsection{The Expansion: Description and Derivation}
\label{sec:20.1}

Wilson~\hyperref[ch20-ref-1]{[1]} hypothesized that the singular part as
$x\to y$ of the product $A(x)B(y)$ of two operators is given by a sum over
other local operators
\begin{equation}
 A(x)B(y)\longrightarrow
 \sum_C F_C^{AB}(x-y)C(y),
 \tag{20.1.1}\label{eq:20.1.1}
\end{equation}
where $F_C^{AB}(x-y)$ are singular c-number functions. Dimensional analysis
suggests that $F_C^{AB}(x-y)$ behaves for $x\to y$ like the power
$d_C-d_A-d_B$ of $x-y$, where $d_O$ is the dimensionality of the operator
$O$ in powers of mass or momentum. Since $d_O$ increases as we add more
fields or derivatives to an operator $O$, the strength of the singularity of
$F_C^{AB}(x-y)$ decreases for operators $C$ of increasing complexity. The
remarkable thing about the operator product expansion is that it is an
\emph{operator} relation; that is, in applying it to any matrix element
$\bra{\beta}A(x)B(y)\ket{\alpha}$, we get the same functions
$F_C^{AB}(x-y)$ for all states $\ket{\alpha}$ and $\ket{\beta}$.

It is the decrease of the singularity in Eq.~\eqref{eq:20.1.1} with
operators $C(y)$ of increasing complexity that makes this expansion useful
in drawing conclusions about the behavior of the product $A(x)B(y)$ for
$x\to y$. The simple power-counting argument above is modified by
renormalization effects; the expansion \eqref{eq:20.1.1} must be formulated
in terms of operators renormalized at some scale $\mu$, and then $\mu$
appears along with $x-y$ in the coefficient function $F_C^{AB}(x-y)$. We
shall see in Section 20.3 that in asymptotically free theories
$F_C^{AB}(x-y)$ does behave like the power $d_C-d_A-d_B$ of $x-y$
suggested by dimensional analysis only up to a power of $\ln(x-y)^2$. Even
in more general theories, it is plausible that the singularities associated
with various operators $C(y)$ will decrease with the complexity of the
operators.

The corresponding statement in momentum space is that for $k\to\infty$,
\begin{equation}
 \int d^4x\,e^{-ik\cdot x}A(x)B(0)
 \longrightarrow \sum_C V_C^{AB}(k)C(0)
 \tag{20.1.2}\label{eq:20.1.2}
\end{equation}
and correspondingly
\begin{equation}
 \int d^4x\,e^{-ik\cdot x}T\{A(x)B(0)\}
 \longrightarrow \sum_C U_C^{AB}(k)C(0),
 \tag{20.1.3}\label{eq:20.1.3}
\end{equation}
where $V_C^{AB}(k)$ and $U_C^{AB}(k)$ are functions of $k^\mu$ that for
large $k$ decrease increasingly rapidly for more and more complicated terms
in the series.

We are going to derive a generalized version of the Wilson expansion, in
which the momenta carried by any number of operators go to infinity together.
For this purpose, let us consider a Greens function for local operators
$A_1(x_1)$, $A_2(x_2)$, etc. whose arguments approach a point $x$, as well
as other local operators $B_1(y_1)$, $B_2(y_2)$, etc. with fixed arguments:
\begin{align}
 &\left\langle
 T\{A_1(x_1),A_2(x_2),\ldots B_1(y_1),B_2(y_2)\ldots\}
 \right\rangle_0
 \notag\\
 &\qquad={}
 \int\left[\prod_{\ell,z}d\phi_\ell(z)\right]
 a_1(x_1)a_2(x_2)\cdots b_1(y_1)b_2(y_2)\cdots
 \exp(iI[\phi]),
 \tag{20.1.4}\label{eq:20.1.4}
\end{align}
where the lower-case letters $a$ and $b$ indicate replacement of the field
operators in the $A$s and $B$s with the c-number fields $\phi$. Now surround
the point $x$ with a ball $B(R)$ of radius $R$ which is much larger than the
separations among the $x_1$, $x_2$, etc. but much smaller than the
separations between $x$, $y_1$, $y_2$, etc. Since the action is local, it may
be written as
\begin{equation}
 I=\int_{z\in B(R)}d^4z\,\mathcal L(z)
  +\int_{z\notin B(R)}d^4z\,\mathcal L(z).
 \tag{20.1.5}\label{eq:20.1.5}
\end{equation}
Eq.~\eqref{eq:20.1.4} may then be put in the form
\begin{align}
 &\left\langle
 T\{A_1(x_1),A_2(x_2),\ldots B_1(y_1),B_2(y_2)\ldots\}
 \right\rangle_0
 \notag\\
 &\quad={}
 \int\left[\prod_{\substack{z\notin B(R)\\ \ell}}d\phi_\ell(z)\right]
 b_1(y_1)b_2(y_2)\cdots
 \exp\left(i\int_{z\notin B(R)}\mathcal L(z)\right)
 \notag\\
 &\qquad{}\cdot
 \int\left[\prod_{\substack{z\in B(R)\\ \ell}}d\phi_\ell(z)\right]
 a_1(x_1)a_2(x_2)\cdots
 \exp\left(i\int_{z\in B(R)}\mathcal L(z)\right),
 \tag{20.1.6}\label{eq:20.1.6}
\end{align}
in which the path integral over the fields inside the ball is constrained by
the boundary condition that the fields merge smoothly at the ball's surface
with the fields outside the ball. Aside from this boundary condition, the
path integral over the fields inside the ball is completely unaffected by
the behavior of the field outside the ball, so the integral over the fields
inside the ball may be expressed in terms of the values and derivatives of
the fields on the surface of the ball, which in turn may be expressed in
terms of the fields and their derivatives extrapolated from outside the ball
to the interior point $x$. If we express this integral as a series in
products\footnote{In order for the coefficients in this series to be finite,
these products must be renormalized by multiplication with suitable infinite
factors. This will be made clearer in the derivation presented in the next
section.} $o(x)$ of the c-number fields and their derivatives extrapolated
to $x$, then the coefficients can only be functions
$U_O^{A_1,A_2,\ldots}(x_1-x,x_2-x,\ldots)$ of the coordinate differences.
Since the points $y_1$, $y_2$, etc. are all far outside the ball $B(R)$, the
exclusion of this ball in the action for the fields outside the ball has no
effect in the limit $R\to0$, so in this limit Eq.~\eqref{eq:20.1.6} becomes
\begin{align}
 &\left\langle
 T\{A_1(x_1),A_2(x_2),\ldots B_1(y_1),B_2(y_2)\ldots\}
 \right\rangle_0
 \longrightarrow
 \int\left[\prod_{\ell,z}d\phi_\ell(z)\right]
 \notag\\
 &\quad{}\cdot b_1(y_1)b_2(y_2)\cdots
 \exp\left(i\int\mathcal L(z)\right)
 \cdot\sum_O
 U_O^{A_1,A_2,\ldots}(x_1-x,x_2-x,\ldots)o(x)
 \notag\\
 &\quad{}=
 \sum_O U_O^{A_1,A_2,\ldots}(x_1-x,x_2-x,\ldots)
 \left\langle
 T\{O(x),B_1(y_1),B_2(y_2)\ldots\}
 \right\rangle_0
 \tag{20.1.7}\label{eq:20.1.7}
\end{align}
for $x_1$, $x_2$, etc. all approaching $x$, where $O(x)$ is the
quantum-mechanical Heisenberg-picture operator corresponding to $o(x)$. In
particular, by Fourier transforming with respect to the $y$ variables and
multiplying with appropriate coefficient functions, this yields
\begin{equation}
 \bra{\beta}T\{A_1(x_1),A_2(x_2),\ldots\}\ket{\alpha}
 \longrightarrow
 \sum_O U_O^{A_1,A_2,\ldots}(x_1-x,x_2-x,\ldots)
 \bra{\beta}O(x)\ket{\alpha}
 \tag{20.1.8}\label{eq:20.1.8}
\end{equation}
for arbitrary states $\ket{\alpha}$ and $\ket{\beta}$. Because this applies
for arbitrary states, it is the operator product expansion in a generalized
version:
\begin{equation}
 T\{A_1(x_1),A_2(x_2),\ldots\}
 \longrightarrow
 \sum_O U_O^{A_1,A_2,\ldots}(x_1-x,x_2-x,\ldots)O(x).
 \tag{20.1.9}\label{eq:20.1.9}
\end{equation}
